\chapter{Variáveis aleatórias discretas}\label{ch:b2:randomvar}

As \omterm{def:b2:randomvar:law}{variáveis aleatórias} organizam os cálculos de probabilidade em torno de
funções, e não de \omterm{def:b2:proba:space}{eventos}. Em espaços \omterm{def:b2:structures:countable}{enumeráveis} a teoria é
movida pelas \omterm{def:b2:series:summable}{famílias somáveis} do~\cref{ch:b2:series}:
a \omterm{def:b2:randomvar:expectation}{esperança} é a soma de uma família indexada pelo \omterm{def:b2:proba:space}{espaço amostral},
e todas as suas propriedades --- linearidade, transferência, fórmula do produto
para variáveis \omterm{def:b2:proba:independence}{independentes} --- são teoremas sobre \omterm{def:b2:series:summable}{famílias
somáveis}. O capítulo demonstra as desigualdades-chave de Markov,
Chebyshev, Cauchy--Schwarz e Jensen, e termina com as \omterm{def:b2:randomvar:law}{leis} clássicas
e a lei fraca dos grandes números, cuja demonstração tem duas linhas
uma vez disponível Chebyshev.

\section{Variáveis aleatórias e suas leis}

\begin{definition}[Variável aleatória discreta; lei]\label{def:b2:randomvar:law}
Seja $(\Omega, \P)$ um \omterm{def:b2:proba:space}{espaço de probabilidade} \omterm{def:b2:structures:countable}{enumerável}. Uma
\emph{variável aleatória}\index{variável aleatória} é uma aplicação $X \colon \Omega \to E$ ($E$ um conjunto
qualquer; variável aleatória \emph{real} quando $E = \R$). Sua \emph{lei}\index{lei!de uma variável aleatória}
(ou \emph{distribuição}\index{distribuição}) é a \omterm{def:b2:proba:space}{medida de probabilidade} $\P_X$ no
\omterm{def:b2:structures:countable}{conjunto enumerável} $X(\Omega)$ definida por
\[
\P_X(\{x\}) = \P(X = x)
= \P\bigl(\{\omega : X(\omega) = x\}\bigr) .
\]
\end{definition}

\begin{example}[As leis clássicas]\label{ex:b2:randomvar:laws}
\leavevmode
\begin{itemize}
\item \emph{Bernoulli} $\mathcal{B}(p)$: $X \in \{0, 1\}$, $\P(X =
1) = p$. Indicadora de um \omterm{def:b2:proba:space}{evento}.
\item \emph{Binomial} $\mathcal{B}(n, p)$: $\P(X = k) =
\binom nk p^k(1-p)^{n-k}$, $0 \leq k \leq n$: número de sucessos
em $n$ provas de Bernoulli \omterm{def:b2:proba:independence}{independentes} (volume do ensino médio; redemonstrada
abaixo via somas de variáveis \omterm{def:b2:proba:independence}{independentes}).
\item \emph{Geométrica} $\mathcal{G}(p)$: $\P(X = k) =
(1-p)^{k-1}p$, $k \in \N^*$: posição do primeiro sucesso
(\cref{ex:b2:proba:geometric}).
\item \emph{Poisson} $\mathcal{P}(\lambda)$: $\P(X = k) =
e^{-\lambda}\frac{\lambda^k}{k!}$, $k \in \N$ --- uma \omterm{def:b2:proba:space}{medida de
probabilidade} pela série exponencial. A lei dos eventos raros
(\cref{ch:b2:genfun}).
\end{itemize}
\end{example}

\begin{remark}[Qual lei modela o quê]
As quatro \omterm{def:b2:randomvar:law}{leis} respondem a quatro perguntas primitivas: Bernoulli,
``aconteceu?''; binomial, ``quantas vezes em $n$
tentativas?''; geométrica, ``quanto tempo até a primeira vez?'';
Poisson, ``quantos \omterm{def:b2:proba:space}{eventos} a uma dada taxa, quando as tentativas são
muitas e individualmente improváveis?''. Reconhecer a pergunta
é nove décimos da modelagem: somas de indicadoras apontam para
a binomial, tempos de espera para a geométrica, contagens de
\omterm{def:b2:proba:space}{eventos} raros para a Poisson --- com a passagem da binomial à
Poisson precisada pela lei dos eventos raros no
\cref{ch:b2:genfun}.
\end{remark}

\begin{proposition}[Ausência de memória da lei
geométrica]\label{prop:b2:randomvar:memoryless}
Se $X \sim \mathcal{G}(p)$, então, para todos $m, n \in \N$:
\[
\P(X > m + n \mid X > m) = \P(X > n) ,
\]
e as \omterm{def:b2:randomvar:law}{leis} geométricas são as únicas \omterm{def:b2:randomvar:law}{leis} em $\N^*$ com essa
propriedade.
\end{proposition}

\begin{proof}
Somando os pesos geométricos, $\P(X > n) = (1-p)^n$. Portanto
\[
\P(X > m + n \mid X > m)
= \frac{\P(X > m + n)}{\P(X > m)}
= \frac{(1-p)^{m+n}}{(1-p)^m} = (1-p)^n = \P(X > n).
\]
Reciprocamente, se $G(n) = \P(X > n)$ satisfaz $G(m + n) =
G(m)G(n)$ com $G(0) = 1$, então $G(n) = G(1)^n$ por indução; $q =
G(1) \in \intco{0}{1}$, e $q = 0$ ou a \omterm{def:b2:randomvar:law}{lei} é
$\mathcal{G}(1 - q)$: $\P(X = k) = G(k-1) - G(k) =
q^{k-1}(1 - q)$.
\end{proof}

\begin{example}[Nenhum número está jamais ``atrasado'']
Lance um dado esperando um seis: o tempo de espera é $X \sim
\mathcal G(1/6)$. A ausência de memória diz que, após $10$
lançamentos infrutíferos, a espera \emph{restante} $X - 10$, dado
$X > 10$, é de novo $\mathcal G(1/6)$: a espera
condicional esperada continua sendo $6$ lançamentos, exatamente como no início.
O dado não se lembra, e nenhum seis está jamais ``atrasado'' --- a
falácia do apostador é a crença de que a \omterm{def:b2:randomvar:law}{lei} condicional
deveria ter mudado. Reciprocamente, a metade de unicidade da proposição
diz que essa indiferença \emph{caracteriza} os tempos de espera
geométricos: todo tempo de espera cuja previsão nunca se atualiza
é geométrico. Filas e tempos de vida reais em geral se atualizam,
que é precisamente como se detecta que eles não são
geométricos.
\end{example}

\section{Esperança}

\begin{definition}[Esperança]\label{def:b2:randomvar:expectation}
Uma \omterm{def:b2:randomvar:law}{variável aleatória} real $X$ em $(\Omega, \P)$ \emph{tem
esperança} se a família
$\bigl(X(\omega)\,\P(\{\omega\})\bigr)_{\omega \in \Omega}$ for
\omterm{def:b2:series:summable}{somável} (\cref{ch:b2:series}); sua \emph{esperança}\index{esperança} é então
\[
\E(X) = \sum_{\omega \in \Omega} X(\omega)\,\P(\{\omega\}) .
\]
\end{definition}

\begin{theorem}[Teorema de transferência]\label{thm:b2:randomvar:transfer}
$X$ tem \omterm{def:b2:randomvar:expectation}{esperança} se e somente se a família $\bigl(x\,\P(X =
x)\bigr)_{x \in X(\Omega)}$ é \omterm{def:b2:series:summable}{somável}, e então
\[
\E(X) = \sum_{x \in X(\Omega)} x\,\P(X = x) .
\]
Mais geralmente, para $f \colon X(\Omega) \to \R$, a variável
$f(X)$ tem \omterm{def:b2:randomvar:expectation}{esperança} se e somente se $\sum_x \abs{f(x)}\,\P(X = x) <
\infty$, e então $\E(f(X)) = \sum_x f(x)\,\P(X = x)$.
\end{theorem}

\begin{proof}
Particione $\Omega$ nos conjuntos de nível $\Omega_x = \{X = x\}$, $x
\in X(\Omega)$. Pelo teorema de soma por pacotes para \omterm{def:b2:series:summable}{famílias somáveis}
(\cref{ch:b2:series}), a família
$(X(\omega)\P(\{\omega\}))_\omega$ é \omterm{def:b2:series:summable}{somável} se e somente se cada pacote o for
(automático: $\sum_{\omega \in \Omega_x}\abs{x}\P(\{\omega\}) =
\abs x\,\P(X = x)$) \emph{e} a família das somas dos pacotes
$\bigl(x\,\P(X = x)\bigr)_x$ for \omterm{def:b2:series:summable}{somável} --- e então as somas
totais coincidem. Para $f(X)$: aplique o enunciado já demonstrado à variável $Y =
f \circ X$, cujos conjuntos de nível são $\{Y = y\} =
\bigsqcup_{x : f(x) = y}\{X = x\}$; uma segunda soma por
pacotes converte $\sum_y y\,\P(Y = y)$ em $\sum_x
f(x)\,\P(X = x)$, agrupando agora os pacotes os valores $x$
por sua imagem $f(x)$, com a \omterm{def:b2:series:summable}{somabilidade} absoluta de uma
família equivalente à da outra.
\end{proof}

\begin{theorem}[Propriedades da esperança]\label{thm:b2:randomvar:linearity}
No conjunto das \omterm{def:b2:randomvar:law}{variáveis aleatórias} com \omterm{def:b2:randomvar:expectation}{esperança}:
\begin{enumerate}
\item (Linearidade) $\E(aX + bY) = a\,\E(X) + b\,\E(Y)$.
\item (Positividade e monotonia) $X \geq 0 \Rightarrow \E(X)
\geq 0$; $X \leq Y \Rightarrow \E(X) \leq \E(Y)$; e
$\abs{\E(X)} \leq \E(\abs X)$.
\item (Dominação) Se $\abs X \leq Z$ e $Z$ tem \omterm{def:b2:randomvar:expectation}{esperança},
então $X$ também tem.
\end{enumerate}
\end{theorem}

\begin{proof}
Todas são propriedades de somas de \omterm{def:b2:series:summable}{famílias somáveis}
(\cref{ch:b2:series}): linearidade da soma, positividade termo a
termo e o critério de dominação para a \omterm{def:b2:series:summable}{somabilidade}. (Note que a
linearidade é imediata na \emph{definição} sobre $\Omega$,
ao passo que seria desajeitada na fórmula de transferência --- um benefício
de definir $\E$ a montante.)
\end{proof}

\begin{example}\label{ex:b2:randomvar:expectations}
$X \sim \mathcal{B}(n, p)$: escrevendo $X = X_1 + \dots + X_n$ como
soma de indicadoras de Bernoulli e usando a linearidade, $\E(X) = np$ ---
sem precisar de coeficientes binomiais. $X \sim \mathcal{G}(p)$: $\E(X) =
\sum_{k\geq1}k(1-p)^{k-1}p = p\cdot\frac{1}{(1 - (1-p))^2} =
\frac1p$, derivando a série geométrica dentro de seu disco
(\cref{ch:b2:powerseries}). $X \sim \mathcal{P}(\lambda)$: $\E(X)
= \sum_{k\geq1}k e^{-\lambda}\frac{\lambda^k}{k!} = \lambda
e^{-\lambda}\sum_{j\geq0}\frac{\lambda^j}{j!} = \lambda$.
\end{example}

\begin{example}[Transferência em ação]\label{ex:b2:randomvar:transferex}
Para $X \sim \mathcal P(\lambda)$, calcule
$\E\bigl(\frac1{1+X}\bigr)$ --- a \omterm{def:b2:randomvar:law}{lei} do próprio $\frac1{1+X}$
é desajeitada, mas a transferência nunca a pede:
\[
\E\Bigl(\frac1{1+X}\Bigr)
= \sum_{k\geq0}\frac{1}{k+1}\,\eu^{-\lambda}
\frac{\lambda^k}{k!}
= \frac{\eu^{-\lambda}}{\lambda}\sum_{k\geq0}
\frac{\lambda^{k+1}}{(k+1)!}
= \frac{\eu^{-\lambda}}{\lambda}\bigl(\eu^\lambda - 1\bigr)
= \frac{1 - \eu^{-\lambda}}{\lambda} .
\]
Duas lições. Computacional: reconhecer uma série exponencial
deslocada é todo o trabalho --- a transferência reduz
\omterm{def:b2:randomvar:expectation}{esperanças} de $f(X)$ a manipulação de séries. Estrutural:
o valor ingênuo por substituição seria $\frac1{1 + \E X} =
\frac1{1 + \lambda}$, e a resposta verdadeira é maior,
\[
\frac{1 - \eu^{-\lambda}}{\lambda} \geq
\frac{1}{1 + \lambda},
\]
exatamente como exige a desigualdade de Jensen para a função \omterm{def:b2:affine:convex}{convexa} $t
\mapsto \frac1{1+t}$. As \omterm{def:b2:randomvar:expectation}{esperanças} de
imagens \omterm{def:b2:affine:convex}{convexas} ficam acima do valor ingênuo por substituição, e a transferência
mais uma verificação de séries torna concreta a desigualdade abstrata.
\end{example}

\begin{theorem}[Independência e produtos]\label{thm:b2:randomvar:product}
As \omterm{def:b2:randomvar:law}{variáveis aleatórias} $X, Y$ são \emph{\omterm{def:b2:proba:independence}{independentes}} se $\P(X = x, Y =
y) = \P(X = x)\P(Y = y)$ para todos $x, y$ --- equivalentemente, os
\omterm{def:b2:proba:space}{eventos} $\{X \in A\}$ e $\{Y \in B\}$ são \omterm{def:b2:proba:independence}{independentes} para todos $A,
B$. Se $X$ e $Y$ são variáveis reais \omterm{def:b2:proba:independence}{independentes} com
\omterm{def:b2:randomvar:expectation}{esperanças}, então $XY$ tem \omterm{def:b2:randomvar:expectation}{esperança} e
\[
\E(XY) = \E(X)\,\E(Y) .
\]
\end{theorem}

\begin{proof}
A equivalência das duas formulações segue somando a identidade
\omterm{def:b2:funcseq:def}{pontual} sobre $(x, y) \in A \times B$
($\sigma$-aditividade duas vezes). Para o produto: a família dupla
$\bigl(xy\,\P(X = x)\P(Y = y)\bigr)_{(x,y)}$ é \omterm{def:b2:series:summable}{somável}, pois, por
Fubini para famílias (\cref{ch:b2:series}),
\[
\sum_{x, y}\abs x \abs y\,\P(X{=}x)\P(Y{=}y)
= \Bigl(\sum_x \abs x \P(X{=}x)\Bigr)
\Bigl(\sum_y \abs y \P(Y{=}y)\Bigr) < \infty ;
\]
pela \omterm{def:b2:proba:independence}{independência} essa família é exatamente $\bigl(xy\,\P(X = x, Y =
y)\bigr)$, cuja soma vale $\E(XY)$ pela transferência aplicada à
variável $(X, Y) \mapsto xy$; Fubini de novo avalia a soma sem
sinais como o produto $\E(X)\E(Y)$.
\end{proof}

\begin{example}[Produtos, com e sem
independência]\label{ex:b2:randomvar:productcontrast}
Lance dois dados honestos. Se $Y$ é o segundo dado (\omterm{def:b2:proba:independence}{independente} do
primeiro), $\E(XY) = \E(X)\E(Y) = 3.5^2 = 12.25$. Se, em vez disso,
$Y = X$ (o ``produto'' de um dado consigo mesmo),
\[
\E(X^2) = \frac{1 + 4 + 9 + 16 + 25 + 36}{6} = \frac{91}{6}
\approx 15.17 \neq 12.25 :
\]
mesmas \omterm{def:b2:randomvar:law}{leis} marginais nos dois cenários, \omterm{def:b2:randomvar:law}{leis} conjuntas diferentes,
\omterm{def:b2:randomvar:expectation}{esperanças} de produto diferentes. A moral, digna de ser gravada:
$\E(XY)$ é um funcional do \emph{par}, não das duas
marginais --- e a diferença $\E(X^2) - \E(X)^2 \approx 2.92$
é, por König--Huygens, precisamente a \omterm{def:b2:randomvar:variance}{variância}
$\frac{35}{12}$ do dado.
\end{example}

\section{Variância, covariância e as desigualdades clássicas}

\begin{definition}[Momentos, variância]\label{def:b2:randomvar:variance}
$X$ tem \emph{momento de ordem 2} se $X^2$ tiver \omterm{def:b2:randomvar:expectation}{esperança}
(e então $X$ também tem, por dominação: $\abs X \leq \frac{1 +
X^2}{2}$). Sua \emph{variância}\index{variância} e seu \emph{desvio
padrão}\index{desvio padrão} são
então
\[
V(X) = \E\bigl((X - \E(X))^2\bigr)
= \E(X^2) - \E(X)^2 ,
\qquad
\sigma(X) = \sqrt{V(X)} ,
\]
(a segunda forma --- a fórmula de \emph{König--Huygens} --- vem de
desenvolver o quadrado e usar a linearidade:
\[
\E\bigl((X - \E X)^2\bigr)
= \E\bigl(X^2 - 2X\,\E X + \E(X)^2\bigr)
= \E(X^2) - 2\,\E(X)^2 + \E(X)^2 ,
\]
o termo do meio usando que $\E X$ é uma constante). Para $X, Y$ com momentos de segunda ordem,
a \emph{covariância}\index{covariância} é
\[
\operatorname{Cov}(X, Y)
= \E\bigl((X - \E X)(Y - \E Y)\bigr)
= \E(XY) - \E(X)\E(Y) .
\]
\end{definition}

\begin{theorem}[Ferramentas da variância]\label{thm:b2:randomvar:variancerules}
Para variáveis com momentos de segunda ordem:
\begin{enumerate}
\item $V(aX + b) = a^2\,V(X)$;
\item $V(X + Y) = V(X) + V(Y) + 2\operatorname{Cov}(X, Y)$ e,
mais geralmente,
\[
V\Bigl(\sum_{i=1}^n X_i\Bigr)
= \sum_{i=1}^n V(X_i)
+ 2\sum_{i < j}\operatorname{Cov}(X_i, X_j) ;
\]
\item se $X, Y$ são \omterm{def:b2:proba:independence}{independentes}, $\operatorname{Cov}(X, Y) = 0$
(a recíproca é falsa), de modo que as \omterm{def:b2:randomvar:variance}{variâncias} de variáveis \omterm{def:b2:proba:independence}{independentes}
se somam.
\end{enumerate}
\end{theorem}

\begin{proof}
\emph{1} e \emph{2} são desenvolvimentos de quadrados mais linearidade;
os produtos $X_iX_j$ têm \omterm{def:b2:randomvar:expectation}{esperança} por Cauchy--Schwarz abaixo
(ou por $\abs{X_iX_j} \leq \frac{X_i^2 + X_j^2}{2}$). \emph{3} é o
\cref{thm:b2:randomvar:product} aplicado às variáveis
centradas. Um contraexemplo padrão para a recíproca: $X$ uniforme
em $\{-1, 0, 1\}$ e $Y = X^2$ são não correlacionadas
($\E(XY) = \E(X^3) = 0 = \E X \cdot \E Y$) mas claramente dependentes.
\end{proof}

\begin{theorem}[Desigualdades de Markov e de
Chebyshev]\label{thm:b2:randomvar:markov}
\leavevmode
\begin{enumerate}
\item (Markov) Se $X \geq 0$ tem \omterm{def:b2:randomvar:expectation}{esperança}, então, para todo $a
> 0$:
\[
\P(X \geq a) \leq \frac{\E(X)}{a} .
\]
\item (Chebyshev) Se $X$ tem momento de segunda ordem, então, para todo
$\varepsilon > 0$:
\[
\P\bigl(\abs{X - \E(X)} \geq \varepsilon\bigr)
\leq \frac{V(X)}{\varepsilon^2} .
\]
\end{enumerate}
\end{theorem}

\begin{proof}
\emph{1.} \omterm{def:b2:funcseq:def}{Pontualmente}, $a\,\mathbf{1}_{X \geq a} \leq X$ (no
\omterm{def:b2:proba:space}{evento} o membro da esquerda vale $a \leq X$; fora dele, $0 \leq X$). Tome
\omterm{def:b2:randomvar:expectation}{esperanças}: $a\,\P(X \geq a) \leq \E(X)$ por monotonia e
$\E(\mathbf{1}_A) = \P(A)$.
\emph{2.} Aplique Markov à variável não negativa $(X - \E
X)^2$ no nível $a = \varepsilon^2$: o \omterm{def:b2:proba:space}{evento} $\{(X - \E X)^2
\geq \varepsilon^2\}$ é exatamente $\{\abs{X - \E X} \geq
\varepsilon\}$.
\end{proof}

\begin{example}[Não correlacionadas mas grudadas
uma na outra]\label{ex:b2:randomvar:sumdiff}
Lance dois dados honestos, $X$ e $Y$ \omterm{def:b2:proba:independence}{independentes}, e ponha $S = X +
Y$, $D = X - Y$. Pela bilinearidade da \omterm{def:b2:randomvar:variance}{covariância},
\[
\operatorname{Cov}(S, D) = V(X) - V(Y) +
\operatorname{Cov}(Y, X) - \operatorname{Cov}(X, Y) = V(X) -
V(Y) = 0 :
\]
soma e diferença são não correlacionadas. \omterm{def:b2:proba:independence}{Independentes}? Certamente
não: $S = 12$ força $D = 0$, enquanto $\P(D = 0) =
\frac16$ incondicionalmente. A correlação só testa a
parte \emph{linear} de uma dependência; aqui a dependência é
carregada pela restrição de que $S$ e $D$ tenham a mesma
paridade, invisível à \omterm{def:b2:randomvar:variance}{covariância}. (Para esse par, a
\omterm{def:b2:randomvar:variance}{covariância} nula precisou de $V(X) = V(Y)$: \omterm{def:b2:randomvar:law}{distribuições} idênticas,
e não \omterm{def:b2:proba:independence}{independência}, é que fizeram o trabalho.)
\end{example}

\begin{example}[Quando Markov é exata]\label{ex:b2:randomvar:markovsharp}
A desigualdade de Markov é uma igualdade precisamente quando nada se
desperdiça na majoração $a\,\mathbf 1_{X\geq a} \leq X$: a
variável deve assumir apenas os valores $0$ e $a$. Concretamente,
se $\P(X = a) = \pi$ e $\P(X = 0) = 1 - \pi$, então $\E(X)
= a\pi$ e
\[
\P(X \geq a) = \pi = \frac{\E(X)}{a} .
\]
Uma leitura realista: numa população em que a riqueza média é
$100$ e a riqueza é $0$ ou $10^6$, a proporção de
milionários é exatamente $10^{-4}$ --- a cota de Markov, atingida
exatamente por desigualdade maximal. Sempre que $X$ se espalha por
valores intermediários a cota é estrita, muitas vezes escandalosamente;
mas, como mostra o caso extremo, nenhuma desigualdade melhor pode ser
extraída apenas da média.
\end{example}

\begin{example}[Chebyshev é ótima --- sem hipóteses
adicionais]\label{ex:b2:randomvar:chebsharp}
Fixe $\varepsilon > 0$, $q \in \intoc01$, e deixe $X$ assumir os
valores $\pm\varepsilon$ com probabilidade $\frac q2$ cada e
$0$ com probabilidade $1 - q$. Então $\E(X) = 0$, $V(X) =
q\varepsilon^2$, e
\[
\P\bigl(\abs{X - \E X} \geq \varepsilon\bigr) = q
= \frac{V(X)}{\varepsilon^2} :
\]
igualdade em Chebyshev. Assim, a desigualdade não pode ser melhorada
usando apenas a \omterm{def:b2:randomvar:variance}{variância} --- o decaimento em $1/\varepsilon^2$ é
o preço exato da informação de segundo momento. Um decaimento mais rápido
exige hipóteses mais fortes: a limitação da variável
compra concentração \emph{exponencial}, como
o~\cref{exo:b2:randomvar:7} antecipa e o problema de fim de semana
deste capítulo desenvolve sistematicamente.
\end{example}

\begin{theorem}[Cauchy--Schwarz e Jensen]\label{thm:b2:randomvar:jensen}
\leavevmode
\begin{enumerate}
\item (Cauchy--Schwarz) Se $X, Y$ têm momentos de segunda ordem, $XY$ tem
\omterm{def:b2:randomvar:expectation}{esperança} e $\E(XY)^2 \leq \E(X^2)\,\E(Y^2)$; consequentemente,
$\operatorname{Cov}(X,Y)^2 \leq V(X)V(Y)$.
\item (Jensen) Se $\varphi \colon I \to \R$ é \omterm{def:b2:affine:convex}{convexa} num
intervalo contendo $X(\Omega)$ e $X$, $\varphi(X)$ têm
\omterm{def:b2:randomvar:expectation}{esperança}, então
\[
\varphi\bigl(\E(X)\bigr) \leq \E\bigl(\varphi(X)\bigr) .
\]
\end{enumerate}
\end{theorem}

\begin{proof}
\emph{1.} \omterm{def:b2:series:summable}{Somabilidade} de $XY$: $\abs{XY} \leq \frac{X^2 + Y^2}2$.
A aplicação $(X, Y) \mapsto \E(XY)$ é uma \omterm{def:b2:quadratic:def}{forma bilinear simétrica}
positiva no espaço das variáveis com momentos de segunda ordem, de modo que a desigualdade
abstrata de Cauchy--Schwarz do~\cref{ch:b2:quadratic} se aplica
(a positividade \emph{semi}definida basta para a desigualdade).
Aplicá-la às variáveis centradas dá a cota da \omterm{def:b2:randomvar:variance}{covariância}.

\emph{2.} Primeiro, $m = \E(X)$ está em $I$: $I$ é um intervalo
contendo todos os valores de $X$, e a \omterm{def:b2:randomvar:expectation}{esperança} é monótona, de modo que $m$
está entre $\inf X(\Omega)$ e $\sup X(\Omega)$. Pelo
teorema da reta de apoio para funções \omterm{def:b2:affine:convex}{convexas}
(\cref{ch:b2:realfun}), existem $\alpha, \beta$ com
$\varphi(t) \geq \alpha t + \beta$ para todo $t \in I$ e
$\varphi(m) = \alpha m + \beta$. Então, \omterm{def:b2:funcseq:def}{pontualmente} em $\Omega$,
$\varphi(X) \geq \alpha X + \beta$; tomando \omterm{def:b2:randomvar:expectation}{esperanças},
\[
\E\bigl(\varphi(X)\bigr) \geq \alpha\,\E(X) + \beta
= \varphi\bigl(\E(X)\bigr). \qedhere
\]
\end{proof}

\begin{example}
Jensen com $\varphi(t) = t^2$ dá $\E(X)^2 \leq \E(X^2)$ ---
a positividade da \omterm{def:b2:randomvar:variance}{variância}; com $\varphi(t) = 1/t$ em
$\intoo{0}{\infty}$: $\frac{1}{\E X} \leq \E\bigl(\frac1X\bigr)$
--- a média harmônica está abaixo da média aritmética, agora em forma
aleatória.
\end{example}

\begin{remark}[Armadilhas comuns]
(i) $\E(XY) = \E(X)\E(Y)$ \emph{exige} \omterm{def:b2:proba:independence}{independência} (ou ao
menos \omterm{def:b2:randomvar:variance}{covariância} nula): tomar $Y = X$ dá $\E(X^2) \neq
\E(X)^2$ sempre que $V(X) > 0$. (ii) Do mesmo modo, $V(X + X) =
4V(X)$, e não $2V(X)$: as \omterm{def:b2:randomvar:variance}{variâncias} só se somam sobre parcelas
\omterm{def:b2:proba:independence}{independentes} (ou não correlacionadas). (iii) $\E(f(X))$ não é
$f(\E(X))$; para $f$ \omterm{def:b2:affine:convex}{convexa}, Jensen até informa a
direção do erro, como no
\cref{ex:b2:randomvar:transferex}. (iv) A existência é uma hipótese
de verdade: para a variável de São Petersburgo $X = 2^K$ com
$\P(K = k) = 2^{-k}$ ($k \geq 1$),
\[
\sum_{k\geq1}2^k\cdot2^{-k} = \sum_{k\geq1}1 = \infty :
\]
$X$ é finita quase certamente e no entanto não tem \omterm{def:b2:randomvar:expectation}{esperança}, e não existe
preço justo de entrada para o jogo. A \omterm{def:b2:series:summable}{somabilidade} na
definição de $\E$ não é pedantismo contábil --- é
onde as caudas pesadas são detectadas. (v) Por fim, o teorema de
transferência precisa da \omterm{def:b2:series:summable}{somabilidade} \emph{absoluta} antes que qualquer
rearranjo da soma sobre os valores seja legítimo
(\cref{ch:b2:series}).
\end{remark}


\begin{example}[Chebyshev em cem
lançamentos]\label{ex:b2:randomvar:hundredtosses}
Para $X \sim \mathcal B(100, \frac12)$: $\E X = 50$, $V(X) =
25$. Chebyshev com $\varepsilon = 6$:
\[
\P(45 \leq X \leq 55) = \P(\abs{X - 50} < 6)
\geq 1 - \frac{25}{36} \approx 0.31 ,
\]
ao passo que a soma binomial exata dá $\approx 0.73$. Os
$31\%$ garantidos estão longe da verdade, mas exigiram
\emph{apenas} a média e a \omterm{def:b2:randomvar:variance}{variância} --- o mesmo certificado
se aplica literalmente a qualquer variável com $\E = 50$, $V = 25$,
por mais exótica que seja, e o~\cref{ex:b2:randomvar:chebsharp} mostra
que alguma tal variável o satura. A universalidade tem um preço;
quando a \omterm{def:b2:randomvar:law}{distribuição} é genuinamente binomial, as ferramentas
exponenciais do problema de fim de semana fecham a maior parte da lacuna.
\end{example}

\begin{example}[A correlação de uma parte com seu
todo]\label{ex:b2:randomvar:partwhole}
Para $X, Y$ \omterm{def:b2:proba:independence}{independentes} e identicamente distribuídas com \omterm{def:b2:randomvar:variance}{variância}
$\sigma^2 > 0$, quão correlacionada é uma parcela com a soma
$S = X + Y$? Calcule
\[
\operatorname{Cov}(X, S) = \operatorname{Cov}(X, X) +
\operatorname{Cov}(X, Y) = \sigma^2 + 0 = \sigma^2,
\qquad V(S) = 2\sigma^2,
\]
de modo que o coeficiente de correlação é
\[
\rho(X, S) = \frac{\operatorname{Cov}(X,
S)}{\sigma(X)\,\sigma(S)}
= \frac{\sigma^2}{\sigma\cdot\sigma\sqrt2}
= \frac{1}{\sqrt2} \approx 0.707 ,
\]
qualquer que seja a \omterm{def:b2:randomvar:law}{lei} comum --- dados, moedas, contagens de Poisson. Com
$n$ parcelas, o mesmo cálculo dá $\rho(X_1, S_n) =
1/\sqrt n$: a influência de cada termo individual sobre o total
se dilui como uma raiz quadrada, que é a sombra correlacional
da escala $\sqrt n$ das flutuações. Cauchy--Schwarz
garante $\abs\rho \leq 1$ sempre; aqui a cota é atingida
exatamente no caso degenerado $n = 1$ e decai previsivelmente
depois.
\end{example}

\begin{example}[MA--MG ponderada a partir de
Jensen]\label{ex:b2:randomvar:amgm}
Seja $Y$ assumindo os valores positivos $a_1, \dots, a_k$ com
probabilidades $\lambda_1, \dots, \lambda_k$. A função
$-\ln$ é \omterm{def:b2:affine:convex}{convexa} em $\intoo0\infty$, de modo que Jensen dá
$-\ln\E(Y) \leq \E(-\ln Y)$, isto é,
\[
a_1^{\lambda_1}a_2^{\lambda_2}\cdots a_k^{\lambda_k}
\;\leq\; \lambda_1a_1 + \lambda_2a_2 + \dots + \lambda_ka_k :
\]
a desigualdade aritmético--geométrica ponderada, com igualdade
se e somente se $Y$ é constante. Pesos iguais $\lambda_i = \frac1k$
recuperam a MA--MG clássica. A probabilidade demonstrou discretamente
um teorema puramente algébrico: escolher uma \omterm{def:b2:randomvar:law}{lei} de probabilidade é
apenas um dispositivo de contabilidade para combinações \omterm{def:b2:affine:convex}{convexas} --- o
ponto de vista baricêntrico do~\cref{ch:b2:affine} mais uma vez, agora
com Jensen como motor.
\end{example}

\section{A lei fraca dos grandes números}

\begin{theorem}[Lei fraca dos grandes números]\label{thm:b2:randomvar:lln}
Sejam $(X_k)_{k \geq 1}$ \omterm{def:b2:randomvar:law}{variáveis aleatórias} \omterm{def:b2:proba:independence}{independentes} duas a duas
com a mesma \omterm{def:b2:randomvar:law}{lei}, admitindo momento de segunda ordem; escreva $m = \E(X_1)$
e $S_n = X_1 + \dots + X_n$. Então, para todo $\varepsilon > 0$:
\[
\P\Bigl(\,\Bigl|\frac{S_n}{n} - m\Bigr| \geq \varepsilon\Bigr)
\;\leq\; \frac{V(X_1)}{n\,\varepsilon^2}
\xrightarrow[n \to \infty]{} 0 .
\]
\end{theorem}

\begin{proof}
Por linearidade, $\E(S_n/n) = m$; pelo
\cref{thm:b2:randomvar:variancerules} (a \omterm{def:b2:proba:independence}{independência} dois a dois
mata as \omterm{def:b2:randomvar:variance}{covariâncias}), $V(S_n) = n\,V(X_1)$, logo $V(S_n/n) =
V(X_1)/n$. A desigualdade de Chebyshev aplicada a $S_n/n$ dá a
cota.
\end{proof}

\begin{remark}
Esse é o teorema que conecta probabilidade a frequência: para
$X_k$ a indicadora de um \omterm{def:b2:proba:space}{evento} $A$ em repetições \omterm{def:b2:proba:independence}{independentes},
$S_n/n$ é a frequência observada de $A$, e a lei dos grandes
números diz que ela se concentra em torno de $\P(A)$ à taxa
$\frac{p(1-p)}{n\varepsilon^2}$. A \omterm{def:b2:randomvar:law}{lei} \emph{forte} ($S_n/n \to
m$ quase certamente) é um teorema do terceiro ano --- sua demonstração para
momentos de quarta ordem está ao alcance, no entanto: veja o~\cref{exo:b2:randomvar:9},
que roda Borel--Cantelli sobre a cota do tipo Chebyshev. A mesma
estimativa de Chebyshev moveu a demonstração por \omterm{thm:b2:funcseq:weierstrass}{polinômios de Bernstein} do
teorema de aproximação de Weierstrass no~\cref{ch:b2:funcseq} --- o
lema de contagem de lá \emph{era} a lei fraca dos grandes números
disfarçada.
\end{remark}

\begin{example}[Colecionando cinquenta cupons]
\label{ex:b2:randomvar:couponnumeric}
O colecionador de cupons do~\cref{exo:b2:randomvar:3} com $n =
50$ brindes distintos: o total esperado vale
\[
\E(T_{50}) = 50\,H_{50} = 50\sum_{k=1}^{50}\frac1k
\approx 50 \times 4.499 \approx 225
\]
caixas --- quatro vezes e meia o palpite ingênuo $50$. O
crescimento harmônico é toda a história: os primeiros $25$ brindes
chegam em cerca de $50\ln2 \approx 35$ caixas, ao passo que o
\emph{último} brinde sozinho custa $50$ caixas em média (uma
espera geométrica de parâmetro $\frac1{50}$). Os problemas
de \omterm{def:b2:metric:complete}{completude} são dominados por seu final de partida, e é por isso que o
\cref{exo:b2:randomvar:12} encontra flutuações da ordem de $n$
--- o tamanho daquela espera geométrica final --- em torno da
média $n\ln n$.
\end{example}

\begin{example}[Quão grande deve ser $n$?]\label{ex:b2:randomvar:llnnumeric}
Para fixar a frequência observada a menos de $\varepsilon = 0.01$ de
$\P(A)$ com confiança $95\%$, a cota de Chebyshev exige
\[
\frac{p(1-p)}{n\varepsilon^2} \leq \frac{1}{4n\varepsilon^2}
\leq 0.05,
\qquad\text{i.e.}\qquad
n \geq \frac{1}{4\cdot0.05\cdot(0.01)^2} = 50\,000 .
\]
A dependência é brutal em $\varepsilon$ (quadrática) e
suave na confiança (linear em $1/\alpha$). Ambas as características
são propriedades da \emph{cota}, não da verdade: as desigualdades
exponenciais do problema de fim de semana baixam o preço da
confiança de $1/\alpha$ para $\ln(1/\alpha)$ --- a mesma
especificação custará ali cerca de $18\,500$ amostras ---
enquanto a escala $1/\varepsilon^2$ é genuína e
inaperfeiçoável. Saber qual parte de uma cota é frouxa é tão
útil quanto a própria cota.
\end{example}

\begin{omfigure}
\begin{tikzpicture}[x=5.5cm, y=1.05cm]
  \draw[->] (-0.05,0) -- (1.1,0) node[below] {$S_n/n$};
  \draw[->] (0,-0.1) -- (0,2.35);
  \draw[dashed] (0.4,0) -- (0.4,2.15);
  \draw[dashed] (0.6,0) -- (0.6,2.15);
  \node[below, font=\scriptsize] at (0.4,0)
    {$m{-}\varepsilon$};
  \node[below, font=\scriptsize] at (0.5,0) {$m$};
  \node[below, font=\scriptsize] at (0.6,0)
    {$m{+}\varepsilon$};
  \draw[omDef, thick] plot[domain=0.06:0.94, samples=80]
    (\x, {0.8*exp(-(\x-0.5)^2/0.0288)});
  \draw[omThm, very thick] plot[domain=0.32:0.68, samples=80]
    (\x, {2.0*exp(-(\x-0.5)^2/0.0032)});
  \node[omDef, font=\small] at (0.16, 0.85) {$n$ pequeno};
  \node[omThm, font=\small] at (0.71, 1.95) {$n$ grande};
\end{tikzpicture}

{\small A lei dos grandes números em imagem: a \omterm{def:b2:randomvar:law}{lei} de
$S_n/n$ (desenhada esquematicamente) mantém seu centro $m$ mas
se estreita à medida que $n$ cresce, de modo que a probabilidade fora da faixa
$\intcc{m-\varepsilon}{m+\varepsilon}$ --- as duas caudas ---
encolhe a zero. Chebyshev majora as caudas por
$V(X_1)/(n\varepsilon^2)$; o problema de fim de semana mostra que elas são
de fato exponencialmente pequenas.}
\end{omfigure}

\begin{remark}[Perspectivas dentro deste volume]
Para a frente, tudo aqui alimenta o~\cref{ch:b2:genfun}: a
\omterm{def:b2:randomvar:expectation}{esperança} $\E(t^X)$ de uma função astuciosa de $X$ empacota
a \omterm{def:b2:randomvar:law}{lei} inteira numa série de potências, os momentos se tornam
derivadas em $1$, e identidades do tipo Wald para somas aleatórias
carregam a teoria dos processos de ramificação; o teorema do produto para
variáveis \omterm{def:b2:proba:independence}{independentes} se torna a multiplicatividade das \omterm{ex:b2:powerseries:fibonacci}{funções
geradoras}. Para trás, a \omterm{def:b2:randomvar:expectation}{esperança} é um \omterm{def:b2:affine:barycenter}{baricentro} com
pesos de probabilidade (\cref{ch:b2:affine}), a desigualdade de Jensen
é a geometria da reta de apoio das funções \omterm{def:b2:affine:convex}{convexas}
(\cref{ch:b2:realfun}) e o método dos momentos exponenciais do
problema de fim de semana deste capítulo é Markov aplicado a
$\eu^{tX}$ --- uma desigualdade, aprimorada por uma boa mudança
de variável, atravessando três capítulos.
\end{remark}

\section{Exercícios}

\begin{exercise}[$\star$]\label{exo:b2:randomvar:1}
Calcule $\E(X)$ e $V(X)$ para $X \sim \mathcal{B}(n, p)$ (via
indicadoras), $X \sim \mathcal{P}(\lambda)$ (mostre que $V(X) =
\lambda$) e $X \sim \mathcal{G}(p)$ (mostre que $V(X) =
\frac{1-p}{p^2}$; use $\E(X(X-1))$ e a segunda derivada da
série geométrica).
\end{exercise}

\begin{exercise}[$\star$]\label{exo:b2:randomvar:2}
Sejam $X \sim \mathcal{P}(\lambda)$ e $Y \sim \mathcal{P}(\mu)$
\omterm{def:b2:proba:independence}{independentes}. Mostre que $X + Y \sim \mathcal{P}(\lambda + \mu)$
(convolução dos pesos; teorema binomial) e que a
\omterm{def:b2:randomvar:law}{lei} condicional de $X$ dado $X + Y = n$ é binomial
$\mathcal{B}\bigl(n, \frac{\lambda}{\lambda + \mu}\bigr)$.
\end{exercise}

\begin{exercise}[$\star$]\label{exo:b2:randomvar:3}
(Colecionador de cupons, \omterm{def:b2:randomvar:expectation}{esperança}) Uma marca de cereal esconde um de $n$
brindes distintos, \omterm{def:b2:funcseq:def}{uniformemente}, em cada caixa. Seja $T_n$ o número de
caixas necessárias para coletar todos os $n$ brindes. Escrevendo $T_n$ como soma de
variáveis geométricas \omterm{def:b2:proba:independence}{independentes} (tempo até ver um brinde \emph{novo} quando
$k$ ainda faltam), mostre que
\[
\E(T_n) = n\sum_{k=1}^{n}\frac{1}{k} \sim n\ln n
\]
(equivalente pela comparação série--integral do
\cref{ch:b2:comparison}).
\end{exercise}

\begin{exercise}[$\star\star$]\label{exo:b2:randomvar:4}
Seja $X \geq 0$ com valores inteiros. Demonstre a \emph{fórmula das caudas}
\[
\E(X) = \sum_{n=1}^{\infty} \P(X \geq n)
\]
(quando um dos membros é finito), escrevendo $X =
\sum_{n\geq1}\mathbf{1}_{X \geq n}$ e trocando as somas
(Fubini para famílias não negativas). Recupere $\E(X) = \frac1p$ para
a \omterm{def:b2:randomvar:law}{lei} geométrica.
\end{exercise}

\begin{exercise}[$\star\star$]\label{exo:b2:randomvar:5}
(A amostragem sem reposição é mais concentrada) Uma urna tem $N$
bolas, $M$ delas brancas. Sorteie $n \leq N$ sem reposição e
seja $X$ a contagem de brancas (\omterm{def:b2:randomvar:law}{lei} \emph{hipergeométrica}). Usando
indicadoras $X = \sum_{i=1}^n Y_i$ com $Y_i$ o $i$-ésimo sorteio:
mostre que cada $Y_i$ é Bernoulli de parâmetro $p = M/N$ (simetria!),
conclua que $\E(X) = np$ exatamente como com reposição e mostre que
$\operatorname{Cov}(Y_i, Y_j) = -\frac{p(1-p)}{N-1} < 0$ para $i
\neq j$, logo $V(X) = np(1-p)\frac{N - n}{N - 1} \leq np(1-p)$.
\end{exercise}

\begin{exercise}[$\star\star$]\label{exo:b2:randomvar:6}
Seja $X$ com momento de segunda ordem. Mostre que $c \mapsto \E\bigl((X -
c)^2\bigr)$ é mínima exatamente em $c = \E(X)$, com mínimo
$V(X)$. Mostre então que $\P(X = \E(X)) = 1$ se e somente se $V(X) =
0$. \emph{(Para o segundo ponto: se $V(X) = 0$, use Chebyshev
com $\varepsilon = 1/n$ e a \omterm{def:b2:metric:continuity}{continuidade} monótona,
\cref{thm:b2:proba:continuity}.)}
\end{exercise}

\begin{exercise}[$\star\star\star$]\label{exo:b2:randomvar:7}
(A concentração bate Markov) Seja $S_n \sim \mathcal{B}(n,
\frac12)$ (número de caras em $n$ lançamentos honestos). Compare as
cotas dadas por Markov ($\P(S_n \geq \frac{3n}{4})$), por Chebyshev
e pelo método exponencial (Chernoff):
\[
\P\Bigl(S_n \geq \frac{3n}4\Bigr)
\leq \E\bigl(e^{tS_n}\bigr)e^{-3nt/4}
= \Bigl(\frac{1 + e^t}{2}\Bigr)^n e^{-3nt/4}
\quad (t > 0),
\]
e otimize $t$ para obter uma cota exponencialmente pequena. \emph{(Em
$t = \ln 3$: majore $\bigl(2\cdot 3^{-3/4}\bigr)^n \approx
(0.877)^n$.)}
\end{exercise}

\begin{exercise}[$\star\star\star$]\label{exo:b2:randomvar:8}
(Weierstrass de novo, probabilisticamente) Seja $f \colon [0,1] \to \R$
\omterm{def:b2:metric:continuity}{contínua} e $S_n \sim \mathcal{B}(n, x)$. Mostre que o
\omterm{thm:b2:funcseq:weierstrass}{polinômio de Bernstein} $B_nf(x) = \sum_{k=0}^n f\bigl(\frac
kn\bigr)\binom nk x^k(1-x)^{n-k}$ vale
$\E\bigl[f\bigl(\frac{S_n}{n}\bigr)\bigr]$ e redemonstre a
estimativa $\abs{B_nf(x) - f(x)} \leq \omega_f(\delta) +
\frac{2\norm f_\infty}{4n\delta^2}$ do~\cref{ch:b2:funcseq} nessa
linguagem probabilística (separe em $\bigl|\frac{S_n}{n} -
x\bigr| \geq \delta$ e use Chebyshev).
\end{exercise}

\begin{exercise}[$\star\star\star$]\label{exo:b2:randomvar:9}
(\omterm{def:b2:randomvar:law}{Lei} forte sob momentos de quarta ordem) Sejam $(X_k)$ \omterm{def:b2:proba:independence}{independentes},
identicamente distribuídas, centradas ($\E X_1 = 0$), com
$\E(X_1^4) < \infty$. Desenvolvendo $\E(S_n^4)$ e contando os
termos sobreviventes (apenas termos em $\E(X_i^4)$ e em $\E(X_i^2X_j^2)$, $i
\neq j$), mostre que $\E(S_n^4) \leq C n^2$ para uma constante $C$. Deduza
$\sum_n \P\bigl(\abs{S_n/n} \geq \varepsilon\bigr) < \infty$
para cada $\varepsilon > 0$ (Markov na ordem 4) e conclua com
Borel--Cantelli (\cref{thm:b2:proba:borelcantelli}) que
$S_n/n \to 0$ quase certamente, numa formulação adequada: o
\omterm{def:b2:proba:space}{evento} $\bigcap_{j}\bigcup_N\bigcap_{n \geq N}\{\abs{S_n/n} <
\frac1j\}$ tem probabilidade $1$.
\end{exercise}

\begin{exercise}[$\star$]\label{exo:b2:randomvar:10}
Lançam-se dois dados honestos; seja $M$ o maior dos dois
resultados. Usando a fórmula das caudas do~\cref{exo:b2:randomvar:4}
(versão finita), mostre que
\[
\E(M) = \sum_{k=1}^{6}\P(M \geq k)
= 6 - \sum_{j=0}^5\Bigl(\frac j6\Bigr)^2 = \frac{161}{36}
\approx 4.47 .
\]
\end{exercise}

\begin{exercise}[$\star\star$]\label{exo:b2:randomvar:11}
Seja $F_n$ o número de pontos fixos de uma permutação uniformemente
aleatória de $\{1, \dots, n\}$ ($n \geq 2$). Escrevendo $F_n =
\sum_i\mathbf 1_{\sigma(i) = i}$, calcule $\E(F_n) = 1$,
$\operatorname{Cov}(\mathbf 1_{\sigma(i)=i}, \mathbf
1_{\sigma(j)=j}) = \frac1{n^2(n-1)}$ para $i \neq j$, e
conclua que $V(F_n) = 1$: em média uma carta fica fixa, com
\omterm{def:b2:randomvar:variance}{variância} exatamente $1$, qualquer que seja $n$.
\end{exercise}

\begin{exercise}[$\star\star\star$]\label{exo:b2:randomvar:12}
(Colecionador de cupons, concentração) No contexto do
\cref{exo:b2:randomvar:3}, mostre que
\[
V(T_n) = \sum_{k=1}^n\frac{1 - k/n}{(k/n)^2}
\leq n^2\sum_{k=1}^n\frac{1}{k^2} \leq \frac{\pi^2}{6}n^2,
\]
usando a \omterm{def:b2:proba:independence}{independência} das etapas geométricas e $V(\mathcal
G(p)) = \frac{1-p}{p^2}$ (\cref{exo:b2:randomvar:1}; o valor
$\pi^2/6$ é \cref{ex:b2:fourier:basel}). Deduza com
Chebyshev que $\dfrac{T_n}{n\ln n} \to 1$ \emph{em
probabilidade}: o tempo total do colecionador é $n\ln n$ a menos de
flutuações da ordem de $n$.
\end{exercise}

\section{Problema: a caixa de ferramentas da concentração, de Markov a
Hoeffding}

\begin{problem}[{Problema de fim de semana --- concentração exponencial
na mão, e quantas pessoas uma pesquisa deve ouvir}]
\label{pb:b2:randomvar:1}
A desigualdade de Markov custa um momento e compra um decaimento em $1/a$;
Chebyshev custa dois momentos e compra $1/\varepsilon^2$ ---
e o~\cref{ex:b2:randomvar:chebsharp} mostra que isso é tudo o que esses
momentos podem comprar. Este problema sobe o resto da escada:
o método exponencial (Chernoff) com sua taxa \emph{exata}
para lançamentos de moeda, a desigualdade de Hoeffding para todas as variáveis
limitadas e a recompensa --- tamanhos de amostra explícitos e honestos
para pesquisas eleitorais, previsões de eleição e teste de moedas.
\index{desigualdades de concentração}\index{desigualdade de
Hoeffding}\index{cota de Chernoff} Em todo o problema, $S_n \sim
\mathcal B(n, p)$ é uma soma de $n$ variáveis de Bernoulli
\omterm{def:b2:proba:independence}{independentes} e $\widehat p_n = S_n/n$ é a frequência empírica.

\medskip
\noindent\textbf{Parte I --- Calibração na moeda honesta.}
Aqui $p = \frac12$ e $a \in \intoo{\frac12}{1}$.
\begin{enumerate}
  \item Markov no nível $an$: mostre que $\P(S_n \geq an) \leq
        \frac1{2a}$, uma cota que não tende sequer a $0$.
        Onde Markov perde tanto?
  \item Chebyshev: usando a simetria da binomial honesta
        em torno de $n/2$, mostre que
        \[
        \P(S_n \geq an) = \tfrac12\,
        \P\bigl(\abs{S_n - \tfrac n2} \geq n(a -
        \tfrac12)\bigr)
        \leq \frac{1}{8n(a - 1/2)^2},
        \]
        isto é, $\frac2n$ em $a = \frac34$: decaimento polinomial,
        enfim.
  \item (Chernoff, nível geral) Calcule
        $\E(\eu^{tS_n}) = \bigl(\frac{1 + \eu^t}2\bigr)^n$ e
        otimize $\P(S_n \geq an) \leq
        \E(\eu^{tS_n})\eu^{-tan}$ em $t > 0$: mostre que o
        $t$ ótimo é $\ln\frac{a}{1-a}$ e que
        \[
        \P(S_n \geq an) \leq \eu^{-n\,I(a)},
        \qquad
        I(a) = \ln 2 + a\ln a + (1-a)\ln(1-a) > 0 .
        \]
        Verifique que $a = \frac34$ recupera a cota
        $\bigl(2\cdot3^{-3/4}\bigr)^n$ do
        \cref{exo:b2:randomvar:7}.
  \item (O expoente é exato) Seja $k = an$ um inteiro.
        A partir do fato de que $\binom nk a^k(1-a)^{n-k}$ é o
        maior dos $n + 1$ termos de uma \omterm{def:b2:randomvar:law}{distribuição} de
        probabilidade, demonstre que $\binom nk \geq
        \frac{\eu^{nH(a)}}{n+1}$ com $H(a) = -a\ln a -
        (1-a)\ln(1-a)$, e deduza a cota inferior correspondente
        \[
        \P(S_n \geq an) \geq \binom{n}{an}2^{-n}
        \geq \frac{\eu^{-n\,I(a)}}{n + 1} .
        \]
  \item Tabule as três cotas em $n = 100$, $a =
        \frac34$: Markov $\frac23$, Chebyshev $0.02$,
        Chernoff $\approx 2.1\cdot10^{-6}$ (o valor verdadeiro é
        $\approx 2.8\cdot10^{-7}$). Moral, numa frase?
\end{enumerate}

\medskip
\noindent\textbf{Parte II --- A desigualdade de Hoeffding.}
\begin{enumerate}[resume]
  \item (Caso Rademacher) Para $\varepsilon = \pm1$ com
        probabilidade $\frac12$ cada, demonstre que
        \[
        \E(\eu^{t\varepsilon}) = \cosh t \leq \eu^{t^2/2}
        \qquad (t \in \R)
        \]
        comparando as duas séries termo a termo ($(2k)! \geq
        2^kk!$).
  \item Deduza, para variáveis de Rademacher \omterm{def:b2:proba:independence}{independentes}
        $\varepsilon_1, \dots, \varepsilon_n$ e todo $s >
        0$:
        \[
        \P\Bigl(\sum_{i=1}^n\varepsilon_i \geq s\Bigr)
        \leq \eu^{-s^2/(2n)} .
        \]
  \item Traduza para moedas honestas ($X_i =
        \frac{1+\varepsilon_i}2$): $\P\bigl(\widehat p_n -
        \tfrac12 \geq \delta\bigr) \leq \eu^{-2n\delta^2}$,
        e a versão bilateral com um fator $2$.
  \item (Lema de Hoeffding) Seja $X \in \intcc01$ com $\E X =
        p$ e $\psi(t) = \ln\E(\eu^{tX})$. Justifique que
        $\psi$ é duas vezes \omterm{def:b2:diffcalc:differential}{diferenciável} com
        \[
        \psi''(t) = \E_t(X^2) - \E_t(X)^2, \qquad
        \E_t(Y) := \frac{\E(Y\eu^{tX})}{\E(\eu^{tX})},
        \]
        uma \emph{\omterm{def:b2:randomvar:variance}{variância}} de uma variável reponderada que ainda
        assume valores em $\intcc01$; majore-a por $\frac14$
        (o argumento de minimalidade do~\cref{exo:b2:randomvar:6}) e
        conclua por Taylor:
        \[
        \E\bigl(\eu^{t(X - p)}\bigr) \leq \eu^{t^2/8} .
        \]
  \item (Desigualdade de Hoeffding) Para $X_i \in
        \intcc01$ \omterm{def:b2:proba:independence}{independentes} com média comum $p$, deduza
        \[
        \P\bigl(\abs{\widehat p_n - p} \geq \delta\bigr)
        \leq 2\,\eu^{-2n\delta^2}
        \qquad (\delta > 0).
        \]
  \item Compare a taxa de Chebyshev
        $\frac{p(1-p)}{n\delta^2}$ com a de Hoeffding
        $2\eu^{-2n\delta^2}$: que hipótese cada uma
        exige, e a partir de qual $n$ (grosseiramente) a
        cota exponencial ganha em $\delta = 0.03$, $p =
        \frac12$?
\end{enumerate}

\medskip
\noindent\textbf{Parte III --- Quantas pessoas uma pesquisa
deve ouvir?} Uma pesquisa consulta $n$ eleitores \omterm{def:b2:proba:independence}{independentes} e uniformemente escolhidos;
cada um responde honestamente; $p$ é a intenção verdadeira e $\widehat p_n$
é o número da pesquisa.
\begin{enumerate}[resume]
  \item Mostre que a pesquisa é precisa a menos de $\pm\delta$ com
        confiança $1 - \alpha$ (isto é,
        $\P(\abs{\widehat p_n - p} \geq \delta) \leq \alpha$)
        assim que
        \[
        n \;\geq\; \frac{\ln(2/\alpha)}{2\,\delta^2} .
        \]
  \item Calcule o $n$ necessário para a especificação padrão ``três
        pontos, noventa e cinco por cento'' ($\delta
        = 0.03$, $\alpha = 0.05$): $n \geq 2050$; e para um
        ponto: $n \geq 18\,445$. Observe --- e explique ---
        o fato notável de que a resposta não envolve
        o tamanho da população.
  \item Refaça a questão 13 com Chebyshev ($V(X_1) = p(1-p)
        \leq \frac14$): $n \geq \frac1{4\alpha\delta^2} =
        5556$ para três pontos. Note que a amostragem
        \emph{sem} reposição só ajuda
        (\cref{exo:b2:randomvar:5}: a \omterm{def:b2:randomvar:variance}{variância} encolhe por
        $\frac{N-n}{N-1}$).
  \item (Prever uma eleição) A intenção verdadeira de um candidato é $p
        = 0.52$. Quantos eleitores devem ser consultados para que
        $\P(\widehat p_n \leq \tfrac12) \leq 0.01$? Mostre que $n
        \geq \frac{\ln 100}{2\cdot(0.02)^2} \approx 5757$ ---
        prever uma disputa apertada custa muito mais que estimar uma
        intenção de voto.
  \item O que a matemática \emph{não} cobre: liste as
        hipóteses de modelagem usadas (amostragem uniforme
        \omterm{def:b2:proba:independence}{independente}, respostas honestas, $p$ fixo) e explique num
        parágrafo curto por que os erros reais de pesquisa são
        dominados pelo \emph{viés} (amostragem não uniforme,
        não resposta), que aumento algum de $n$ reduz.
\end{enumerate}

\medskip
\noindent\textbf{Parte IV --- Mais fino e mais barato.}
\begin{enumerate}[resume]
  \item (Mediana das médias: decaimento exponencial a partir de dois momentos)
        Divida um orçamento de $km$ amostras em $k$ grupos
        independentes de $m$; sejam $\widehat p^{(1)}, \dots, \widehat
        p^{(k)}$ as médias dos grupos e $M$ a mediana delas.
        Escolha $m$ de modo que cada grupo satisfaça
        $\P(\abs{\widehat p^{(i)} - p} \geq \delta) \leq
        \frac18$ (Chebyshev: $m \geq \frac2{\delta^2}$
        basta). Mostre que, se $\abs{M - p} \geq \delta$,
        então ao menos $k/2$ grupos erram, e deduza
        \[
        \P(\abs{M - p} \geq \delta)
        \leq \binom{k}{\lceil k/2\rceil}\Bigl(\frac18
        \Bigr)^{k/2}
        \leq 2^k\cdot 8^{-k/2} = 2^{-k/2} :
        \]
        concentração exponencial sem usar nada além de
        \omterm{def:b2:randomvar:variance}{variâncias}.
  \item (Paley--Zygmund) Para $X \geq 0$ com momento de segunda ordem,
        demonstre que $\P(X > 0) \geq \dfrac{\E(X)^2}{\E(X^2)}$
        \emph{(Cauchy--Schwarz em $X\mathbf 1_{X>0}$)}: a
        ferramenta na direção inversa --- os momentos também podem forçar
        \omterm{def:b2:proba:space}{eventos} a acontecer.
  \item (Pinsker leve) Mostre que $I(a) \geq 2\bigl(a -
        \tfrac12\bigr)^2$ em $\intoo{\frac12}1$ \emph{(a
        diferença se anula até a segunda ordem em $\frac12$ e
        sua segunda derivada é $\frac1{a(1-a)} - 4 \geq
        0$)}: o expoente exato de Chernoff sempre bate o
        quadrático de Hoeffding.
  \item Desenvolva $I\bigl(\tfrac12 + \delta\bigr) = 2\delta^2 +
        O(\delta^4)$ e combine com a questão 4: para desvios
        pequenos, o expoente de Hoeffding $2n\delta^2$ é
        assintoticamente \emph{exato} --- método algum pode batê-lo
        por mais que fatores polinomiais.
  \item Monte a tabela da caixa de ferramentas: para Markov, Chebyshev, a
        cota de quarto momento do~\cref{exo:b2:randomvar:9},
        Hoeffding e Chernoff com expoente $I$, enuncie em
        uma linha cada: hipótese exigida, decaimento obtido
        e a questão deste problema em que ela foi
        mais fina.
\end{enumerate}

\medskip
\noindent\textbf{Parte V --- Dividendos.}
\begin{enumerate}[resume]
  \item (Testando uma moeda) Uma moeda é honesta ou é viciada com
        $p = 0.55$. Você a lança $n$ vezes e declara
        ``viciada'' quando $\widehat p_n > 0.525$. Mostre que
        ambas as probabilidades de erro são no máximo
        $\eu^{-2n(0.025)^2}$ e que $n \geq 3685$ lançamentos
        garantem as duas abaixo de $1\%$.
  \item (\omterm{def:b2:proba:space}{Eventos} raros pedem uma cota atenta à \omterm{def:b2:randomvar:variance}{variância}) Seja $p =
        0.01$ e tome a especificação relativa $\delta =
        p/2 = 0.005$, $\alpha = 0.05$. Compare os tamanhos
        de amostra exigidos por Hoeffding ($n \approx 74\,000$) e
        por Chebyshev com a \omterm{def:b2:randomvar:variance}{variância} verdadeira $p(1-p)$ ($n
        \approx 7920$): a cota exponencial cega à \omterm{def:b2:randomvar:variance}{variância}
        perde para o humilde momento de segunda ordem. Enuncie a moral
        e diga de onde virá a ferramenta que falta (uma cota
        exponencial atenta à \omterm{def:b2:randomvar:variance}{variância}; a aproximação de Poisson do
        \cref{ch:b2:genfun}).
  \item (\omterm{def:b2:randomvar:law}{Lei} forte para moedas) A partir do $\sum_n
        2\eu^{-2n\delta^2} < \infty$ e de Borel--Cantelli
        (\cref{thm:b2:proba:borelcantelli}), demonstre que
        $\widehat p_n \to p$ quase certamente para lançamentos de moeda
        \omterm{def:b2:proba:independence}{independentes}: formule o \omterm{def:b2:proba:space}{evento} quase certo como
        $\bigcap_j\bigcup_N\bigcap_{n\geq N}
        \{\abs{\widehat p_n - p} < \tfrac1j\}$, como no
        \cref{exo:b2:randomvar:9}, e conclua. (A limitação
        substitui o momento de quarta ordem usado lá.)
  \item Síntese. Em cinco frases: o que cada degrau da
        escada (momentos um, dois, quatro; exponencial limitado;
        expoente exato) custa e compra; por que consultar $2050$
        pessoas basta para um país de qualquer tamanho; e qual
        dessas cotas o volume do terceiro ano de graduação afinará nas
        constantes exatas do teorema central do limite.
\end{enumerate}
\end{problem}
